\section{Peano arithmetic}\label{sec:peano_arithmetic}

\begin{concept}\label{con:arithmetic}
  In a broad sense, the study of numbers is named \term[ru=арифметика (\cite[48]{БашмаковаЮшкевич1951ПроисхождениеСистемСчисления}), en=arithmetic (\cite[3]{Stillwell2010MathHistory})]{arithmetic} after the Greek word \enquote{αριθμος} (\enquote{number}).

  \begin{itemize}
    \thmitem{con:arithmetic/elementary} Arithmetic originates as the study of the so-called \term[ru=арифметические операции (\cite[101]{Проскуряков1951ТеоретическиеОсновыАрифметики}), en=arithmetic operator (\cite[131]{Knuth1997ArtVol1})]{arithmetic operations} --- addition, subtraction, multiplication and division --- and is mostly concerned with integers.

    A historical overview of this \enquote{elementary} arithmetic can be found in \cite{Karpinski1925HistoryOfArithmetic}. A notable book of historical interest on the topic is \bycite{Fibonacci2002LiberAbaci}.

    We will present systematized results in \fullref{sec:positional_number_systems} and \fullref{sec:digit_based_operations}. Some modern considerations stemming from computer science are presented in \fullref{sec:fixed_width_arithmetic}.

    \thmitem{con:arithmetic/algebraic} Ever since antiquity, the study of integers goes way beyond the arithmetic operations. \hyperref[def:prime_number]{Prime numbers} and \hyperref[def:polynomial_equation]{Diophantine equations} have notably led to some of the most abstract mathematics in existence.

    The abstract study of numbers is colloquially called \term[ru=теория чисел (\cite[7]{Бухштаб1966ТеорияЧисел}), en=theory of numbers (\cite[1]{Apostol1976AnalyticNumberTheory})]{number theory}, although some influential books like \bycite{Gauss1801DisquisitionesArithmeticae} and \bycite{Serre1973Arithmetic} call it \enquote{arithmetic}.

    \thmitem{con:arithmetic/logical} The development of mathematical logic in the 19th century is tied to the formalization of arithmetic. Using the formalisms from \fullref{ch:mathematical_logic}, we will present the theory \logic{PA} --- Peano arithmetic --- in \cref{def:peano_arithmetic}.

    The relation of this theory to the work of Peano and others is briefly discussed in \cref{rem:peano_arithmetic}.
  \end{itemize}
\end{concept}

\paragraph{Peano axioms}

\begin{definition}\label{def:peano_arithmetic}\mcite[def. 2.7.7; 150]{VanDalen2004LogicAndStructure}
  \term[ru=арифметика Пеано (\cite[167]{Герасимов2014Вычислимость}), en=Peano arithmetic (\cite[def. 4.1.1(ii)]{Hinman2005Logic})]{Peano arithmetic}, abbreviated as \logic{PA}, is a \hyperref[def:fol_theory]{theory} of \hyperref[sec:first_order_logic]{first-order predicate logic} for describing \hyperref[def:natural_numbers]{natural numbers} and \hyperref[con:arithmetic/elementary]{arithmetic operations} on them.

  We will present two theories in addition to \logic{PA} --- the theory \( \logic{PA}_2 \) of \hyperref[def:second_order_logic]{monadic second-order logic}, which replaces the induction schema \eqref{eq:def:peano_arithmetic/theory/PA3} with \eqref{eq:def:peano_arithmetic/theory/PA3/sol}, and \term[en=Robinson arithmetic (\cite[def. 4.1.1(i)]{Hinman2005Logic})]{Robinson arithmetic} \logic{Q}, which omits the induction schema altogether. We discuss the relation of these theories to Peano's work in \cref{rem:peano_arithmetic}.

  \begin{thmenum}
    \thmitem{def:peano_arithmetic/signature} Consider the \hyperref[def:fol_signature]{first-order signature} \( \Sigma_{\logic{PA}} \) containing the following symbols:
    \begin{thmenum}
      \thmitem{def:peano_arithmetic/signature/zero} A constant \( \syn0 \) for representing zero.

      \thmitem{def:peano_arithmetic/signature/succ} A unary \hyperref[def:fol_signature/function]{functional symbol} \( \synS \) using \hyperref[def:fol_signature/notation]{condensed notation}, called the \term[en=successor (\cite[87]{VanDalen2004LogicAndStructure})]{successor operation}.

      \thmitem{def:peano_arithmetic/signature/plus} An \hyperref[def:fol_signature/notation]{infix} binary functional symbol \( \synplus \) for denoting addition.

      \thmitem{def:peano_arithmetic/signature/times} Another infix binary functional symbol \( \synmult \) for denoting multiplication.
    \end{thmenum}

    For \( \logic{PA}_2 \), we can translate this signature using \fullref{alg:fol_to_hol/signature}. Technically, to avoid collisions with the syntax for \hyperref[def:lambda_term_schema]{\( \muplambda \)-term placeholders}, we need to replace the symbol \( \synS \) with another symbol like \( \synS^+ \).

    \thmitem{def:peano_arithmetic/theory} Over \( \Sigma_{\logic{PA}} \), consider the theory \( \logic{PA} \) with the following axioms:
    \begin{thmenum}
      \thmitem{def:peano_arithmetic/theory/PA1} Zero is not the successor of any natural number:
      \begin{equation}\label{eq:def:peano_arithmetic/theory/PA1}\tag{\logic{PA1}}
        \synneg \qexists \synn (\synS \synn \syneq \syn0).
      \end{equation}

      This ensures \hyperref[def:well_founded_relation]{well-foundedness} of the natural numbers.

      \thmitem{def:peano_arithmetic/theory/PA2} The successor function is \hyperref[thm:function_invertibility_categorical/nonempty_left_invertible]{injective}:
      \begin{equation}\label{eq:def:peano_arithmetic/theory/PA2}\tag{\logic{PA2}}
        \qforall \synn \qforall \synm (\synS \synn \syneq \synS \synm \synimplies \synn \syneq \synm).
      \end{equation}

      \thmitem{def:peano_arithmetic/theory/PA3} The \term[ru=аксиомы индукции (\cite[168]{Герасимов2014Вычислимость}), en=induction schema (\cite[def. 2.7.7]{VanDalen2004LogicAndStructure})]{axiom schema of induction} provides a powerful proof method:
      \begin{equation}\label{eq:def:peano_arithmetic/theory/PA3}\tag{\logic{PA3}}
        \parens[\big]
          {
            \underbrace{\syn\varphi[\synx \synsubst \syn0]}_{\T{base case}}
            \synwedge
            \qforall \synn \parens[\big]
              {
                \overbrace
                  {
                    \underbrace{ \syn\varphi[\synx \synsubst \synn] }_{\mathclap{\substack{\T{inductive} \\ \T{hypothesis}}}}
                    \synimplies
                    \underbrace{ \syn\varphi[\synx \synsubst \synS \synn] }_{\mathclap{\substack{\T{inductive step} \\ \T{conclusion}}}}
                  }^{\T{inductive step}}
              }
          }
        \synimplies
        \underbrace{ \qforall \synn \syn\varphi[\synx \synsubst \synn] }_{\T{conclusion}}.
      \end{equation}

      The instance of \( \varphi \) may have free variables, so we must take the \hyperref[def:fol_quantifier_closure]{universal closure} of each concrete \hyperref[con:schemas_and_instances]{instance}.

      Note that \eqref{eq:def:peano_arithmetic/theory/PA3} is a formula schema in the sense of \cref{def:fol_formula_schema}; it is an adaptation of the following second-order axiom\fnote{We discuss the emulation of second-order logic via schemas in \cref{rem:simulating_second_order_logic}}:
      \begin{equation}\label{eq:def:peano_arithmetic/theory/PA3/sol}\tag{\( \logic{PA3}_\logic{2} \)}
        \qforall {\synp^{\syn\iota \synimplies \syn\omicron}}
        \parens[\big]
          {
            \parens[\big]
              {
                \synp \syn0
                \synwedge
                \qforall {\synn^{\syn\iota}} \parens
                  {
                    \synp \synn
                    \synimplies
                    \synp \synS^+ \synn
                  }
              }
            \synimplies
            \qforall {\synn^{\syn\iota}} \synp \synn
          }.
      \end{equation}

      The schema is stated as an \hyperref[def:inference_rule]{inference rule} in \cref{thm:peano_arithmetic_induction_rule}. A metatheoretic restatement of is given in \fullref{thm:induction_on_natural_numbers}, while a variation of the principle is given in \fullref{thm:strong_induction_on_natural_numbers}.

      We discuss induction more generally in \cref{con:induction}.

      \thmitem{def:peano_arithmetic/theory/plus} The next two axioms (or rather, their universal closures) inductively define the desired behavior of addition:
      \begin{align}
        \synn \synplus \syn0       &\syneq \synn,                       \label{eq:def:peano_arithmetic/theory/PA4}\tag{\logic{PA4}} \\
        \synn \synplus \synS \synm &\syneq \synS(\synn \synplus \synm). \label{eq:def:peano_arithmetic/theory/PA5}\tag{\logic{PA5}}
      \end{align}

      \thmitem{def:peano_arithmetic/theory/times} The (universal closures of) final two axioms are for multiplication:
      \begin{align}
        \synn \synmult \syn0       &\syneq \syn0,                           \label{eq:def:peano_arithmetic/theory/PA6}\tag{\logic{PA6}} \\
        \synn \synmult \synS \synm &\syneq (\synn \synmult \synm) + \synn. \label{eq:def:peano_arithmetic/theory/PA7}\tag{\logic{PA7}}
      \end{align}
    \end{thmenum}
  \end{thmenum}
\end{definition}
\begin{comments}
  \item The phrase \enquote{model of Peano arithmetic} may refer to either \hyperref[def:fol_structure]{first-order logic structures} that satisfy \logic{PA} or to \hyperref[def:hol_structure]{higher-order logic structures} that satisfy \( \logic{PA}_2 \). On a technical level the difference is significant, however, when dealing with standard models in the sense of \cref{def:peano_arithmetic_standard_model}, the two concepts are equivalent.

  Nonstandard models, on the other hand, can be rather weird --- see \cref{ex:peano_arithmetic_nonstandard_models}.

  \item These axioms are formalized in \cite{notebook:code}, in the modules
  \begin{CodeRefDisplay}
    \ManualCodeRef{logic.theories.arithmetic.axioms} \\
    \ManualCodeRef{lambda_.hol.theories.arithmetic.axioms}
  \end{CodeRefDisplay}
\end{comments}

\begin{remark}\label{rem:peano_arithmetic}
  We have presented several \hyperref[def:logical_theory]{logical theories} for \hyperref[def:natural_numbers]{natural number} \hyperref[con:arithmetic]{arithmetic} in \cref{def:peano_arithmetic}. We discuss here how they are related to the work of Giuseppe Peano.

  \begin{thmenum}
    \thmitem{rem:peano_arithmetic/sets} \incite*[1]{Peano1889ArithmeticesPrincipia} not only presents a formal theory of arithmetic, but, more importantly, he formulates an early variant of \hyperref[con:predicate_logic]{predicate logic} based on \hyperref[def:set_theory_signature]{set theory}.

    Peano's axioms are most faithfully represented by what \incite[70]{Enderton1977SetTheory} and \incite[53]{Mendelson2008NumberSystems} call \enquote{Peano systems} --- triples \( (N, e, S) \) such that
    \begin{itemize}
      \item \( S \) is an \hi{injective} endofunction on \( N \) (injectivity subsumes \eqref{eq:def:peano_arithmetic/theory/PA1}).

      \item \( e \) is an element of \( N \) that is not in the range of \( S \) (the latter condition subsumes \eqref{eq:def:peano_arithmetic/theory/PA2}).

      \item If \( A \) contains \( e \) and is closed under \( S \), then \( A \) is a subset of \( N \) (restatement of \eqref{eq:def:peano_arithmetic/theory/PA3/sol}).
    \end{itemize}

    Peano systems resemble Dedekind's earlier construction from \cite[def. 73]{Dedekind2007NatureAndMeaningOfNumbers}. Peano himself refines his logical system (and theory) in \cite[27]{Peano1908FormulatioMathematico}.

    Neither of the aforementioned authors consider addition and multiplication to be \hyperref[con:primitive_notion]{primitive notions} --- this later turns out to be a necessity in the first-order theory.

    \thmitem{rem:peano_arithmetic/zero} Both \incite[1]{Peano1889ArithmeticesPrincipia} and \incite[def. 73]{Dedekind2007NatureAndMeaningOfNumbers} consider \( 1 \) --- not \( 0 \) --- to be the first natural number. Peano's more refined theory in \cite[27]{Peano1908FormulatioMathematico} is adjusted to that \( 0 \) is the first natural number.

    This dichotomy is also present in other author's work --- when defining Peano systems, Enderton starts with \( 0 \), while Mendelson starts with \( 1 \).

    We find it convenient for \( 0 \) to be a natural number because this allows \( \BbbN \) to become a \hyperref[def:monoid]{monoid} and hence a \hyperref[def:semiring]{semiring}.

    On the other hand, we find it more convenient to start counting from \( 1 \).

    In the presence of ambiguity, we will distinguish between \enquote{nonnegative integers} and \enquote{positive integers} --- these are defined formally in \cref{def:integer_ordering}.

    \thmitem{rem:peano_arithmetic/successor} Although Peano introduces a dedicated successor operation, he denotes it by \( a + 1 \) in \cite[1]{Peano1889ArithmeticesPrincipia} and by \( a+ \) in \cite[27]{Peano1908FormulatioMathematico}.

    Using a dedicated symbol for this operation is a modern practice. That the successor can be expressed via addition can be proven easily --- see \cref{thm:peano_arithmetic_successor_and_addition}.

    \thmitem{rem:peano_arithmetic/fol_variants} Modern formulations of first-order Peano arithmetic mostly agree with each other, up to a precise choice of formalisms:
    \begin{itemize}
      \item Our list of axioms coincides with \bycite[def. 2.7.7]{VanDalen2004LogicAndStructure} and \bycite[\S 5.2.5]{Mimram2020ProgramEqualsProof}.

      \item \incite[105]{КолмогоровДрагалин2006Логика} and \incite[167]{Герасимов2014Вычислимость} extend our list with axioms for equality. We will not need them because we handle equality as a primordial logical primitive.

      \item \incite*[def. 4.1.1(ii)]{Hinman2005Logic} extends our signature with the strict inequality \( \synless \) and adds corresponding axioms. We instead present \( \synleq \) and related predicates via \hyperref[def:fol_definitional_extension]{definitional extensions} in \cref{def:peano_arithmetic_ordering}.
    \end{itemize}
  \end{thmenum}
\end{remark}

\begin{theorem}[Induction on natural numbers]\label{thm:induction_on_natural_numbers}
  The \hyperref[def:peano_arithmetic]{Peano arithmetic} induction axioms \eqref{eq:def:peano_arithmetic/theory/PA3} and \eqref{eq:def:peano_arithmetic/theory/PA3/sol} can be restated semantically.

  Fix a \hyperref[def:fol_semantics/model]{model} \( \mscrX = (X, I) \) of Peano arithmetic. For clarity, denote \( I(\syn0) \) by \( 0 \) and \( I(\synS) \) by \( S \).

  Fix also a \hyperref[def:boolean_function]{Boolean-valued function} \( p: X \to \set{ \semtop, \sembot } \). For models of the second-order theory \( \logic{PA}_2 \), \( p \) can be arbitrary, but for the first-order theory \logic{PA}, \( p \) must be the \hyperref[def:fol_parameterized_formula_denotation]{parametrized denotation} of a formula with one free variable.

  In order for \( p(n) \) to hold for every \( n \) in \( X \), the following is sufficient:
  \begin{thmenum}
    \thmitem{thm:induction_on_natural_numbers/base} Base case: \( p(0) \) must hold. For other possible base cases, see \cref{rem:natural_number_induction_base}.

    \thmitem{thm:induction_on_natural_numbers/step} Inductive step: \( p(n) \) must entail \( p(S(n)) \).
  \end{thmenum}
\end{theorem}
\begin{comments}
  \item A useful variation of this theorem is \fullref{thm:strong_induction_on_natural_numbers}.

  See \cref{con:induction} for a general discussion of induction and recursion.
\end{comments}
\begin{proof}
  \SubProof{Proof of first-order case} Fix a formula \( \varphi \) and a variable assignment \( v \) into \( X \). Suppose that \( \Bracks{\varphi}_\mscrX^{v_{x \to 0}} = \semtop \) and, for any \( n \) in \( X \), \( \Bracks{\varphi}_\mscrX^{v_{x \to n}} \leq \Bracks{\varphi}_\mscrX^{v_{x \to S(n)}} \).

  \Cref{thm:fol_formula_semantics_of_assignment_substitution} implies that
  \begin{equation*}
    \Bracks{\varphi[x \mapsto \syn0]}_\mscrX^v = \semtop
  \end{equation*}
  and similarly, for any \( n \),
  \begin{equation*}
    \Bracks{\varphi[x \mapsto n]}_\mscrX^{v_{x \to n}} \leq \Bracks{\varphi[x \mapsto \synS n]}_\mscrX^{v_{x \to n}},
  \end{equation*}
  which due to \cref{thm:def:heyting_algebra/leq} is equivalent to
  \begin{equation*}
    \Bracks{\varphi[x \mapsto n] \synimplies \varphi[x \mapsto \synS n]}_\mscrX^{v_{x \to n}} = \semtop.
  \end{equation*}

  Then, by the definition \eqref{eq:alg:fol_formula_denotation/forall} of the denotation of universal formulas,
  \begin{equation*}
    \Bracks{\qforall n (\varphi[x \mapsto n] \synimplies \varphi[x \mapsto \synS n])}_\mscrX^v = \semtop.
  \end{equation*}

  Both premises of \eqref{eq:def:peano_arithmetic/theory/PA3} hold, so \fullref{thm:propositional_semantic_mp} implies that its conclusion also does.

  \SubProof{Proof of second-order case} We proceed similarly, but without the schema instantiation mechanism. We replace \( \varphi \) with applications of the variable \( \synp^{\syn\iota \synimplies \syn\omicron} \), and we use assignments \( v \) for which \( v(\synp) = p \).
\end{proof}

\begin{proposition}\label{thm:peano_arithmetic_induction_rule}
  The induction axiom \eqref{eq:def:peano_arithmetic/theory/PA3} can be restated as the following \hyperref[def:natural_deduction_rule]{natural deduction rule}:
  \begin{equation*}\taglabel[\logic{PA3}]{inf:thm:peano_arithmetic_induction_rule}
    \begin{prooftree}
      \hypo{ \syn\varphi[\synx \synsubst \syn0] }
      \hypo{ \qforall \synn (\syn\varphi[\synx \synsubst \synn] \synimplies \syn\varphi[\synx \synsubst \synS \synn]) }
      \infer2[\ref{inf:thm:peano_arithmetic_induction_rule}]{ \qforall \synn \syn\varphi[\synx \synsubst \synn] }
    \end{prooftree}
  \end{equation*}

  This rule is \hyperref[con:inference_rule_admissibility]{admissible} in \hyperref[def:peano_arithmetic]{Peano arithmetic}.
\end{proposition}
\begin{proof}
  \SubProofOf{thm:peano_arithmetic_induction_rule}
  \begin{equation*}
    \begin{prooftree}
      \hypo{ [\eqref{eq:def:peano_arithmetic/theory/PA3}]^{a_1} }
      \infer1[\ref{inf:def:fol_natural_deduction/forall/elim}*]{ \synanon }

      \hypo{ [\varphi[x \mapsto \syn0]]^u }
      \hypo{ [\qforall n (\varphi_n \synimplies \varphi_{\synS n})]^w }

      \infer2[\ref{inf:def:propositional_natural_deduction/and/intro}]{ [\varphi[x \mapsto \syn0] \synwedge \qforall n (\varphi_n \synimplies \varphi_{\synS n})]^v }

      \infer2[\ref{inf:def:propositional_natural_deduction/imp/elim}]{ \qforall n \varphi_n }
    \end{prooftree}
  \end{equation*}

  Note that, if \( \varphi \) has free variables, we apply \ref{inf:def:fol_natural_deduction/forall/elim} as many times as needed to the axiom instance.
\end{proof}

\begin{proposition}\label{thm:peano_arithmetic_predecessor_existence}
  In \hyperref[def:peano_arithmetic]{Peano arithmetic}, every number except zero has a predecessor. More precisely, the following is provable:
  \begin{equation}\label{eq:thm:peano_arithmetic_predecessor_existence}
    \qforall \synn ((\synn \syneq \syn0) \synvee (\qexists \synm \synn \syneq \synS \synm)).
  \end{equation}
\end{proposition}
\begin{comments}
  \item An interesting feature of this proof is that it uses induction, but only as a way to traverse the natural numbers --- the inductive hypothesis is never used.
\end{comments}
\begin{proof}

  Let
  \begin{equation*}
    \varphi \coloneqq (\synx \syneq \syn0) \synvee (\qexists \synm \synx \syneq \synS \synm).
  \end{equation*}

  Then
  \begin{equation*}
    \begin{prooftree}
      \infer0[\ref{inf:def:fol_natural_deduction/equality/intro}]{ \syn0 \syneq \syn0 }
      \infer1[\ref{inf:def:propositional_natural_deduction/or/intro_left}]{ \varphi[\synx \mapsto \syn0] }

      \infer0[\ref{inf:def:fol_natural_deduction/equality/intro}]{ (\synS \synn \syneq \synS \synm)[\synm \mapsto \synn] }
      \infer1[\ref{inf:def:fol_natural_deduction/exists/intro}]{ \qexists \synm (\synS \synn \syneq \synS \synm) }
      \infer1[\ref{inf:def:propositional_natural_deduction/or/intro_right}]{ \varphi[\synx \mapsto \synS \synn] }
      \infer1[\ref{inf:def:propositional_natural_deduction/imp/intro}]{ \varphi[\synx \mapsto \synn] \synimplies \varphi[\synx \mapsto \synS \synn] }
      \infer1[\ref{inf:def:fol_natural_deduction/forall/intro}]{ \qforall \synn \parens[\big]{ \varphi[\synx \mapsto \synn] \synimplies \varphi[\synx \mapsto \synS \synn] } }

      \infer2[\ref{inf:thm:peano_arithmetic_induction_rule}]{ \qforall \synn \varphi[\synx \mapsto \synn] }
    \end{prooftree}
  \end{equation*}
\end{proof}

\begin{proposition}\label{thm:peano_arithmetic_successor_and_addition}
  In \hyperref[def:peano_arithmetic]{Peano arithmetic}, we can express the successor operation by adding \( 1 = S(0) \). More precisely, the following is provable:
  \begin{equation}\label{eq:thm:peano_arithmetic_successor_and_addition}
    \qforall \synn (\synS \synn \syneq \synn \synplus \synS \syn0).
  \end{equation}
\end{proposition}
\begin{comments}
  \item When not dealing with formal arithmetic, we will avoid the successor operation altogether.
  \item We extend this to \hyperref[def:ordinal]{ordinal numbers} in \cref{rem:ordinal_successor_via_addition}.
\end{comments}
\begin{proof}
  \begin{equation*}
    \begin{prooftree}
      \hypo{ [\eqref{eq:def:peano_arithmetic/theory/PA5}]^{a_1} }
      \infer1[\ref{inf:def:fol_natural_deduction/forall/elim}*]{ \synn \synplus \synS \syn0 \syneq \synS(\synn \synplus \syn0) }
      \infer1[\ref{inf:thm:fol_natural_deduction_equality/symm}]{ \synS(\synn \synplus \syn0) \syneq \synn \synplus \synS \syn0 }

      \hypo{ [\eqref{eq:def:peano_arithmetic/theory/PA4}]^{a_2} }
      \infer1[\ref{inf:def:fol_natural_deduction/forall/elim}]{ \synn \synplus \syn0 \syneq \synn }

      \infer2[\ref{inf:def:fol_natural_deduction/equality/elim}]{ \synS \synn \syneq \synn \synplus \synS \syn0 }
      \infer1[\ref{inf:def:fol_natural_deduction/forall/intro}]{ \qforall \synn (\synS \synn \syneq \synn \synplus \synS \syn0) }
    \end{prooftree}
  \end{equation*}
\end{proof}

\begin{remark}\label{rem:natural_number_induction_base}
  When using \fullref{thm:induction_on_natural_numbers}, it is sometimes necessary to consider base cases other than \( n = 0 \).

  There are examples all across mathematics --- \cref{thm:def:rational_numbers_ordering/power_monotone}, \cref{thm:real_number_positive_power}, \cref{thm:prime_ideal_is_radical}, \cref{thm:def:lambda_term_alpha_equivalence/free}, etc. --- but we will demonstrate a toy example here, cited verbatim from \cite{Cohen1961NatureOfMathematicalProofs}:
  \begin{displayquote}
    \textbf{\textsc{lemma I.}} \textit{All horses are the same color} (by induction).

    \quad \textit{Proof:} It is obvious that one horse is the same color. Let us assume the proposition, \( P(k) \), that \( k \) horses are the same color and show this to imply that \( k + 1 \) horses are the same color. Given the set of \( k + 1 \) horses, we remove one horse; then the remaining \( k \) horses are the same color, by hypothesis. We remove another horse and replace the first; the \( k \) horses, by hypothesis, are again the same color. We repeat this until by exhaustion the \( k + 1 \) sets of \( k \) horses each have shown to be the same color. It follows that, since every horse is the same color as every other horse, \( P(k) \) entails \( P(k + 1) \). But since we have shown \( P(1) \) to be true, \( P \) is true for all succeeding values of \( k \); i.e. all horses are the same color.
  \end{displayquote}

  The proof of the inductive step relies on ambiguity. We must pick horses \( a \) and \( b \) from a set \( H \) of \( k + 1 \) horses. By the inductive hypothesis, all horses in \( H \setminus \set{ a } \) and in \( H \setminus \set{ b } \) will have the same color. Any horse \( c \) in \( H \setminus \set{ a, b } \) must then have the same color as \( a \), and similarly for \( c \) and \( b \), so it remains for \( a \) and \( b \) to have the same color.

  This reasoning only works in \( H \) has at least three elements, in which case the inductive hypothesis is \( P(2) \). This requires both \( P(1) \) and \( P(2) \) to be proven explicitly.

  To analyze this from the perspective of \hyperref[def:peano_arithmetic]{Peano arithmetic}, let us consider some extension \( \Sigma^+ \) of the first-order signature \( \Sigma_{\logic{PA}} \) with a unary predicate \( \synP \), which is to be interpreted as \( P \).

  Let us now consider the instance of \eqref{eq:def:peano_arithmetic/theory/PA3} with \( \varphi = \synP(\synx) \):
  \begin{equation*}
    \parens[\big]
      {
        \underbrace{\synP(\syn0)}_{\mathclap{\T{base case}}}
        \synwedge
        \qforall \synn \parens[\big]
          {
            \overbrace
              {
                \underbrace{ \synP(\synn) }_{\mathclap{\substack{\T{inductive} \\ \T{hypothesis}}}}
                \synimplies
                \underbrace{ \synP(\synS \synn) }_{\mathclap{\substack{\T{ind. step} \\ \T{conclusion}}}}
              }^{\T{inductive step}}
          }
      }
    \synimplies
    \underbrace{ \qforall \synn \synP(\synn) }_{\T{conclusion}}
  \end{equation*}

  There are several things to note here:
  \begin{itemize}
    \item Cohen considers \( k = 1 \) to be the base case rather than \( k = 0 \). This general dichotomy on whether \( 0 \) is considered a natural number is discussed in \cref{rem:peano_arithmetic/zero}. What matters here is that \( P(0) \) holds vacuously, so in this case the issue is irrelevant. If this is not the case, we can instead use
    \begin{equation*}
      \varphi = (\synx \syneq \syn0) \synvee \synP(\synx).
    \end{equation*}

    \item An additional base case does not change the induction axiom itself, but complicates the proof trees. For example, a straightforward proof of the inductive step has the following form:
    \begin{equation*}
      \begin{prooftree}
        \hypo{ [\synP(\synn)]^i }
        \infer[dashed]1{ \synP(\synS \synn) }
      \end{prooftree}
    \end{equation*}

    An additional base case for \( k = 1 \) requires the following structure instead:
    \begin{equation*}
      \begin{prooftree}
        \infer0[\ref{inf:thm:propositional_admissible_rules/lem}]{ (\synn \syneq \syn0) \synvee \synneg (\synn \syneq \syn0) }

        \hypo{ [\synn \syneq \syn0]^u }
        \infer[dashed]1{ \synP(\synS \synn) }

        \hypo{ [\synneg (\synn \syneq \syn0)]^v }
        \hypo{ [\synP(\synn)]^i }
        \infer[dashed]2{ \synP(\synS \synn) }

        \infer[left label={\( u, v \)}]3[\ref{inf:def:propositional_natural_deduction/or/elim}]{ \synP(\synS \synn) }
      \end{prooftree}
    \end{equation*}

    Adding a second or third base case requires similar extensions.
  \end{itemize}
\end{remark}

\paragraph{Models of Peano arithmetic}

\begin{definition}\label{def:peano_arithmetic_term_model}\mimprovised
  \Fullref{thm:omega_is_model_of_pa} implies that the \hyperref[thm:smallest_inductive_set_existence]{smallest inductive set} \( \omega \) acts as a \hyperref[def:fol_semantics/model]{model} of \hyperref[def:peano_arithmetic]{Peano arithmetic} that is standard in the sense of \cref{def:peano_arithmetic_standard_model}, i.e. a first-order structure that satisfies the second-order axioms.

  We will build another first-order model \( \mscrN = (N, I) \) based on \( \omega \), with the advantage of the elements of \( N \) themselves being \hyperref[def:fol_term]{logical terms} that can be interpreted in any other structure. We will call it the \term{standard term model} of \logic{PA} since it resembles the term model construction from \cref{def:fol_term_model}, and it is standard in the sense of \cref{def:peano_arithmetic_standard_model}.

  First, we use \fullref{thm:omega_recursion} to construct the \hyperref[def:sequence]{sequence}
  \begin{equation*}
    \tau_k \coloneqq \begin{cases}
      \syn0,        &k = \varnothing, \\
      \synS \tau_m, &k = \op{sc}(m).
    \end{cases}
  \end{equation*}

  Let \( N \coloneqq \set{ \tau_k \given k \in \omega } \). As we will discuss in \cref{rem:def:peano_arithmetic_term_model}, thus defined, \( N \) faithfully represents the informal linguistic definition
  \begin{equation*}
    N \coloneqq \set{ \syn0, \synS \syn0, \synS \synS \syn0, \synS \synS \synS \syn0, \ldots }.
  \end{equation*}

  As a function, \( \tau: \omega \to N \) is \hyperref[def:function_invertibility/bijective]{bijective}. This allows us to identify elements of \( \omega \) and \( N \), and to unambiguously define an interpretation \( I \) for \( N \) based on \( \omega \).

  \Fullref{alg:fol_to_hol} provides us with a \hyperref[def:hol_structure]{standard second-order structure} \( \mscrN_2 \) --- a \hyperref[def:hol_semantics/model]{model} of \( \logic{PA}_2 \).
\end{definition}
\begin{defproof}
  We must show that, as a function, \( \tau \) is really bijective. It is clearly surjective, so we must show injectivity. We will use \fullref{thm:omega_induction} on \( i \) to show that \( \tau_i = \tau_j \) implies \( i = j \).
  \begin{itemize}
    \item If \( i = \varnothing \), then \( \tau_i = \syn0 = \tau_j \) is the unique string in \( N \) that starts with \( \syn0 \), and hence \( j = \varnothing \).

    \item Suppose that \( i = \op{sc}(k) \), where the inductive hypothesis holds for \( k \). By definition, \( \tau_i = \synS \tau_k = \tau_j \), so there must exist some \( m \) such that \( j = \op{sc}(m) \). In particular, \( \tau_j = \synS \tau_m \), where \( \tau_k = \tau_m \).

    By the inductive hypothesis on \( k \), we conclude that \( k = m \), hence also \( i = j \).
  \end{itemize}
\end{defproof}

\begin{remark}\label{rem:def:peano_arithmetic_term_model}
  Consider the standard term model \( \mscrN = (N, I) \) of \hyperref[def:peano_arithmetic]{Peano arithmetic} from \cref{def:peano_arithmetic_term_model}.

  We have defined \( N \) as the set \( \set{ \tau_k \given k \in \omega } \), where
  \begin{equation*}
    \tau_k \coloneqq \begin{cases}
      \syn0,        &k = \varnothing, \\
      \synS \tau_m, &k = \op{sc}(m),
    \end{cases}
  \end{equation*}
  with the intention to formalize the informal linguistic definition
  \begin{equation*}
    N \coloneqq \set{ \syn0, \synS \syn0, \synS \synS \syn0, \synS \synS \synS \syn0, \ldots }.
  \end{equation*}

  The recursive definition depends on precise constructions from \fullref{sec:naive_set_theory} based on \hyperref[def:inductive_set]{inductive sets}. On the other hand, the informal definition is based on \fullref{sec:formal_languages}, which in turn depend on primordial \hyperref[con:metalogic]{metatheoretic} natural numbers. So, our task is to show that the concepts of natural numbers in the object theory and metatheory coincide.

  For the sake of definiteness, suppose that the linguistic definition amounts to the set \( C \) of all \hyperref[def:fol_closed_term]{closed first-order term} over the signature consisting only of \( \syn0 \) and \( \synS \).

  By definition, \( N \) is a subset of \( C \). We will use the metatheoretic \fullref{thm:induction_on_natural_numbers} on the \hyperref[def:formal_language/string_length]{string length} of the term \( \sigma \) from \( C \) to show that \( \tau \) belongs to \( N \).
  \begin{itemize}
    \item The shortest term is \( \syn0 \) of length \( 1 \), which equals \( \tau_\varnothing \).

    \item Suppose that the inductive hypothesis holds for \( n \) and suppose that \( \sigma \) has length \( n + 1 \). Then \( \sigma = a \rho \) for some symbol \( a \), which can only be \( \synS \).

    Then \( \rho \) is a well-formed term of length \( n \), so the inductive hypothesis implies that \( \rho = \tau_k \) for some \( k \in \omega \). It follows that \( \rho = \synS \tau_k = \tau_m \), where \( m = \op{sc}(k) \).
  \end{itemize}
\end{remark}

\begin{proposition}\label{thm:peano_arithmetic_standard_segment}
  Via \fullref{alg:fol_term_denotation}, the \hyperref[def:peano_arithmetic_term_model]{standard term model} of \hyperref[def:peano_arithmetic]{Peano arithmetic} (\logic{PA}) \hyperref[def:fol_embedding]{embeds} into any \hyperref[def:fol_semantics/model]{first-order model} of \logic{PA}.
\end{proposition}
\begin{comments}
  \item In other words, every first-order model \( \mscrX \) has a submodel isomorphic to \( \mscrN \). \incite[164]{Hinman2005Logic} calls this submodel the \enquote{standard part} of \( \mscrX \).
\end{comments}
\begin{proof}
  Let \( \mscrN = (N, I) \) be the standard term model. Fix any model \( \mscrY = (Y, J) \) of \logic{PA} and consider the map
  \begin{equation*}
    \begin{aligned}
      &\Phi: N \to Y, \\
      &\Phi(\tau) \coloneqq \Bracks{\tau}_\mscrY.
    \end{aligned}
  \end{equation*}

  We must show that \( \Phi \) is an embedding, i.e. an injective homomorphism (it is automatically a strong homomorphism since the signature contains no predicate symbols).

  \SubProofOf[def:function_invertibility/injective/equality]{injectivity} We will use \fullref{thm:induction_on_natural_numbers} on the term \( \tau_1 \) to show that \( \Phi(\tau_1) = \Phi(\tau_2) \) implies \( \tau_1 = \tau_2 \). This property can be formalized an abstract \hyperref[def:boolean_function]{Boolean-valued function} defined in the metalanguage, so we utilize the fact that the translation \( \mscrN_2 \) satisfies the second-order axiom \eqref{eq:def:peano_arithmetic/theory/PA3/sol}.

  \begin{itemize}
    \item If \( \tau_1 = \syn0 \), then necessarily \( \tau_2 = \syn0 \).

    Indeed, if \( \tau_2 = \synS \sigma_2 \) for some element \( \sigma_1 \) of \( N \), then
    \begin{equation*}
      \Bracks{\syn0}_\mscrY = \Phi(\tau_1) = \Phi(\tau_2) = \Bracks{\synS \sigma_2}_\mscrY.
    \end{equation*}

    \Cref{thm:fol_equality_characterization} implies that \( \mscrY \) satisfies \( \syn0 \syneq \synS \sigma_2 \), which contradicts the assumption that \( \mscrY \) satisfies \eqref{eq:def:peano_arithmetic/theory/PA1}.

    \item Otherwise, \( \tau_1 = \synS \sigma_1 \), where the inductive hypothesis holds for \( \sigma_1 \).

    If \( \tau_2 = \syn0 \), we can similarly derive a contradiction, so \( \tau_2 = \synS \sigma_2 \) for some \( \sigma_2 \). Then
    \begin{equation*}
      \Bracks{\synS \sigma_1}_\mscrY = \Phi(\tau_1) = \Phi(\tau_2) = \Bracks{\synS \sigma_2}_\mscrY.
    \end{equation*}

    From \cref{thm:fol_equality_characterization} it follows that \( \mscrY \) satisfies the formula \( \synS \sigma_1 \syneq \synS \sigma_2 \). Since \( \mscrY \) also satisfies \eqref{eq:def:peano_arithmetic/theory/PA2}, via \fullref{thm:propositional_semantic_mp} we conclude that \( \mscrY \) satisfies \( \sigma_1 \syneq \sigma_2 \). Then \cref{thm:fol_equality_characterization} implies that \( \Bracks{\sigma_1}_\mscrY = \Bracks{\sigma_2}_\mscrY \).

    The inductive hypothesis on \( \sigma_1 \) implies that \( \sigma_1 = \sigma_2 \). Therefore,
    \begin{equation*}
      \synS \sigma_1 = \tau_1 = \tau_2 = \synS \sigma_2.
    \end{equation*}
  \end{itemize}

  \SubProof{Proof that \( \Phi \) is a homomorphism} We must show that \( \Phi \) respects the functions from \( \Sigma_{\logic{PA}} \).

  \SubProof*{Proof that \( \Phi \) respects \( \syn0 \)} By definition, \( I(\syn0) = \syn0 \), so
  \begin{equation*}
    \Phi(I(\syn0)) = \Phi(\syn0) = \Bracks{\syn0}_\mscrY = J(\syn0).
  \end{equation*}

  \SubProof*{Proof that \( \Phi \) respects \( \synS \)} Similarly,
  \begin{equation*}
    \Phi(I(\synS)(\tau)) = \Phi(\synS \tau) = \Bracks{\synS \tau}_\mscrY = J(\synS)(\Bracks{\tau}_\mscrY) = J(\synS)(\Phi(\tau)).
  \end{equation*}

  \SubProof*{Proof that \( \Phi \) respects \( \synplus \)} We must use \fullref{thm:induction_on_natural_numbers} on \( \sigma \) to show
  \begin{equation*}
    \Phi(I(\synplus)(\tau, \sigma)) = J(\synplus)(\Phi(\tau), \Phi(\sigma)).
  \end{equation*}

  \begin{itemize}
    \item Suppose that \( \sigma = \syn0 \). \Cref{thm:fol_formula_semantics_of_compatible_substitutions} implies that, for any variable assignment \( v \) into \( N \),
    \begin{equation*}
      \Bracks{\tau \synplus \syn0 \syneq \tau}_\mscrN^v
      =
      \Bracks{\underbrace{\synx \synplus \syn0 \syneq \synx}_{\eqref{eq:def:peano_arithmetic/theory/PA4}}}_\mscrN^{v_{\synx \mapsto \tau}}
      =
      \semtop.
    \end{equation*}

    Then \cref{thm:fol_equality_characterization} implies that
    \begin{equation*}
      I(\synplus)(\tau, \sigma) = \Bracks{\tau \synplus \syn0}_\mscrN = \Bracks{\tau}_\mscrN = \tau.
    \end{equation*}

    Similarly, we conclude that
    \begin{equation*}
      J(\synplus)(\Phi(\tau), \Phi(\sigma)) = \Bracks{\tau \synplus \syn0}_\mscrY = \Bracks{\tau}_\mscrY = \Phi(\tau).
    \end{equation*}

    Therefore,
    \begin{equation*}
      \Phi(I(\synplus)(\tau, \sigma)) = J(\synplus)(\Phi(\tau), \Phi(\sigma)).
    \end{equation*}

    \item Similarly, if \( \sigma = \synS \rho \), where the inductive hypothesis holds for \( \rho \), we have
    \begin{equation*}
      I(\synplus)(\tau, \sigma) = \Bracks{\tau \synplus \synS \rho}_\mscrN = \Bracks{\synS (\tau \synplus \rho)}_\mscrN = I(\synS)\parens[\big]{ I(\synplus)(\tau, \rho) }
    \end{equation*}
    and
    \begin{equation*}
      J(\synplus)(\Phi(\tau), \Phi(\sigma)) = \cdots = J(\synS)\parens[\big]{ J(\synplus)\parens[\big]{ \Phi(\tau), \Phi(\rho) } }.
    \end{equation*}

    By the inductive hypothesis on \( \rho \), \( \Phi(I(\synplus)(\tau, \rho)) = J(\synplus)(\Phi(\tau), \Phi(\rho)) \), therefore
    \begin{balign*}
      \Phi(I(\synplus)(\tau, \sigma))
      &=
      \Phi\parens[\big]{ I(\synS)\parens[\big]{ I(\synplus)(\tau, \rho) } }
      = \\ &=
      J(\synS)\parens[\big]{ \Phi\parens[\big]{ I(\synplus)(\tau, \rho) } }
      = \\ &=
      J(\synS)\parens[\big]{ J(\synplus)(\Phi(\tau), \Phi(\rho)) }
      = \\ &=
      J(\synplus)(\Phi(\tau), \Phi(\sigma)),
    \end{balign*}
    as desired.
  \end{itemize}

  \SubProof*{Proof that \( \Phi \) respects \( \synmult \)} Similar.
\end{proof}

\begin{definition}\label{def:peano_arithmetic_standard_model}\mimprovised
  We say that a \hyperref[def:fol_semantics/model]{first-order model} of \hyperref[def:peano_arithmetic]{Peano arithmetic} is \term{standard} if the evaluation embedding of the \hyperref[def:peano_arithmetic_term_model]{standard term model} from \cref{thm:peano_arithmetic_standard_segment} is an \hyperref[def:fol_isomorphism]{isomorphism}.

  Nonstandard models have \enquote{excessive} elements compared to the standard models listed in \cref{thm:peano_arithmetic_standard_models}. We call these excessive elements \term[en=non-standard elements (\cite[def. 2.5.10]{Hinman2005Logic})]{nonstandard natural numbers}.
\end{definition}
\begin{comments}
  \item Due to \cref{thm:peano_arithmetic_categoricity}, this definition aligns with \cref{rem:simulating_second_order_logic}, where a model of first-order \hyperref[def:peano_arithmetic]{Peano arithmetic} (\logic{PA}) is considered standard if its translation via \fullref{alg:fol_to_hol} is a \hyperref[def:hol_semantics/model]{model} of the second-order theory \( \logic{PA}_2 \).

  \item This definition is loosely based on the discussion of \hyperref[def:hol_theory/categorical]{second-order categoricity} in \cite[\S 5]{Farmer2008STTVirtues}.

  From the perspective of first-order logic, nonstandard models of arithmetic are discussed in \cite[\S 2.5]{Hinman2005Logic}, \cite[\S 3.2]{VanDalen2004LogicAndStructure} and \cite[\S 3.3.2]{Герасимов2014Вычислимость}.

  \item The importance of standard models is highlighted in \cref{ex:peano_arithmetic_fol_induction}; see \cref{ex:peano_arithmetic_nonstandard_models} for examples of nonstandard models.
\end{comments}

\begin{proposition}\label{thm:peano_arithmetic_categoricity}
  \hyperref[def:peano_arithmetic]{Second-order Peano arithmetic} is \hyperref[def:hol_theory/categorical]{categorical} with respect to \hyperref[def:hol_semantics/standard]{standard semantics}: every model of \( \logic{PA}_2 \) is \hyperref[def:hol_isomorphism]{isomorphic} to the (second-order) \hyperref[def:peano_arithmetic_term_model]{standard term model} \( \mscrN_2 \).
\end{proposition}
\begin{proof}
  Fix a second-order model \( \mscrY = (Y, J) \) of \( \logic{PA}_2 \), where \( Y = \seq{ Y_\tau }_{\tau \in \op*{\BbbQ_{\Sigma_{\logic{PA}}}}} \) is a \hyperref[def:hol_frame]{standard frame}.

  We will show that the map
  \begin{equation*}
    \begin{aligned}
      &\Phi: N \to Y_{\syn\iota}, \\
      &\Phi(\tau) \coloneqq \Bracks{\tau}_\mscrY
    \end{aligned}
  \end{equation*}
  is an isomorphism. By an adaptation of the arguments in \cref{thm:peano_arithmetic_standard_segment}, we can conclude that \( \Phi \) is an embedding. We only need to show that it is surjective.

  We use \fullref{thm:induction_on_natural_numbers} on the element \( n \) of \( Y_{\syn\iota} \) to prove the existence of a term \( \tau \) of \( N \) such that \( \Phi(\tau) = n \). As in the proof of \cref{thm:peano_arithmetic_standard_segment}, we use induction on an abstract Boolean-valued function defined in the metatheory, so it is important for \( \mscrY \) to satisfy the second-order axiom \eqref{eq:def:peano_arithmetic/theory/PA3/sol}.

  \begin{itemize}
    \item For the base case \( n = J(\syn0) \), naturally \( \tau \coloneqq \syn0 \).
    \item Suppose that \( n = J(\synS)(m) \), where the inductive hypothesis holds for \( m \).

    Then there exists some term \( \sigma \) from \( N \) such that \( \Phi(\sigma) = m \). It remains to define \( \tau \coloneqq \synS \sigma \).
  \end{itemize}
\end{proof}

\begin{example}\label{ex:peano_arithmetic_fol_induction}
  We list examples where the first-order induction axiom schema \eqref{eq:def:peano_arithmetic/theory/PA3} fails to capture the predicate needed for \fullref{thm:induction_on_natural_numbers}:
  \begin{thmenum}
    \thmitem{ex:peano_arithmetic_fol_induction/embedding} In \cref{thm:peano_arithmetic_standard_segment}, when proving that the \hyperref[def:peano_arithmetic_term_model]{standard term model} of \hyperref[def:peano_arithmetic]{Peano arithmetic} (\logic{PA}) embeds into any first-order model of \logic{PA}, we use induction several times without encoding the inductive hypothesis as a first-order formula. In fact, no such formulas exist because these hypotheses use formalisms outside of the scope of the object theory.

    \thmitem{ex:peano_arithmetic_fol_induction/categoricity} In \cref{thm:peano_arithmetic_categoricity}, when showing that the aforementioned embedding is surjective for any standard second-order model \( \mscrX \) of \( \logic{PA}_2 \), we utilize induction on \( \mscrX \), but again provide no first-order formula for similar reasons.
  \end{thmenum}
\end{example}

\begin{proposition}\label{thm:peano_arithmetic_standard_models}
  The following are \hyperref[def:peano_arithmetic_standard_model]{standard models} of \hyperref[def:peano_arithmetic]{Peano arithmetic}:
  \begin{itemize}
    \thmitem{thm:peano_arithmetic_standard_models/omega} The \hyperref[thm:smallest_inductive_set_existence]{smallest inductive set} \( \omega \).

    \thmitem{thm:peano_arithmetic_standard_models/terms} The \hyperref[def:peano_arithmetic_term_model]{standard term model} \( \mscrN \).

    \thmitem{thm:peano_arithmetic_standard_models/initial} Any \hyperref[def:universal_objects/initial]{initial object} in the \hyperref[def:semiring/category]{category of semirings}.
  \end{itemize}
\end{proposition}
\begin{comments}
  \item This proposition merely highlights different possibilities, which may or may not depend on an already established notion of natural numbers.
\end{comments}
\begin{proof}
  \SubProofOf{thm:peano_arithmetic_standard_models/omega} \Cref{thm:omega_is_model_of_pa} states that \( \omega \) is a model of second-order Peano arithmetic.

  \SubProofOf{thm:peano_arithmetic_standard_models/terms} The standard term models is defined to match both the universe and the interpretation of \( \omega \), so naturally it is also a model.

  \SubProofOf{thm:peano_arithmetic_standard_models/initial} \Cref{thm:category_of_semirings_properties/initial} implies that any standard model of \( \logic{PA} \) is an initial object in the category of unital semirings, and \cref{thm:def:universal_objects/initial} implies that any other initial object is isomorphic to it.
\end{proof}

\begin{example}\label{ex:peano_arithmetic_nonstandard_models}
  We list examples of \hyperref[def:peano_arithmetic_standard_model]{nonstandard models} of \hyperref[def:peano_arithmetic]{Peano arithmetic}:
  \begin{thmenum}
    \thmitem{ex:peano_arithmetic_nonstandard_models/lowenheim_skolem} We have shown in \cref{thm:peano_arithmetic_standard_models} that there are standard models.

    Standard models are necessarily countable in the metatheory; indeed, the \hyperref[def:peano_arithmetic_term_model]{standard term model} \( \mscrN \) consists of strings over a finite alphabet.

    Using \fullref{thm:upward_lowenheim_skolem_theorem}, we can construct an \hyperref[def:elementary_embedding]{elementary extension} \( \mscrM \) of \( \mscrN \) of any desired infinite cardinality, say \hyperref[def:aleph_hierarchy]{\( \aleph_1 \)}.

    Then \( \mscrM \) contains uncountably many nonstandard natural numbers not present in \( N \). Yet, \( \mscrN \) and \( \mscrM \) satisfy the same first-order sentences.

    \thmitem{ex:peano_arithmetic_nonstandard_models/compactness} A simpler example of a nonstandard model of \logic{PA} can be constructed by using \fullref{thm:fol_semantic_compactness} directly.

    Let \( \Sigma^+ \) be the extension of \( \Sigma_{\logic{PA}} \) by a single individual constant, \( \sync \). Let \( \logic{PA}^+ \) be the extension of \logic{PA} with the following axioms:
    \begin{align*}
      \synneg (\sync &\syneq \syn0), \\
      \synneg (\sync &\syneq \synS \syn0), \\
      \synneg (\sync &\syneq \synS \synS \syn0), \\
                     &\vdots
    \end{align*}

    Any finite subset of \( \logic{PA}^+ \) contains only finitely many such axioms, so the standard term model \( \mscrN \) contains at least one closed term not subject to them. Therefore, \( \logic{PA}^+ \) is \hyperref[def:finitely_satisfiable_set_of_sentences]{finitely satisfiable}.

    \Fullref{thm:fol_semantic_compactness} implies that \( \logic{PA}^+ \) is \hyperref[def:fol_semantics/satisfaction]{satisfiable}. Then there exists a model \( \mscrM \) of \( \logic{PA}^+ \), which is incidentally also a model of \( \logic{PA} \).

    \Cref{thm:peano_arithmetic_standard_segment} implies that \( \mscrN \) embeds into \( \mscrM \). Because of the new axioms imposed on \( \mscrM \), we can use \fullref{thm:induction_on_natural_numbers} to show that the denotation of \( \sync \) in \( \mscrM \) differs from the denotation of any term in \( \mscrN \).

    Therefore, \( \mscrM \) contains a nonstandard natural number.
  \end{thmenum}
\end{example}

\paragraph{Incompleteness}

\begin{theorem}[Peano arithmetic incompleteness]\label{thm:peano_arithmetic_incompleteness}\mcite[thm. 5.5.37]{Герасимов2014Вычислимость}
  \hyperref[def:peano_arithmetic]{First-order Peano arithmetic} is \hyperref[def:fol_theory/complete]{incomplete}.
\end{theorem}
\begin{comments}
  \item We discuss the history of this theorem in \cref{rem:thm:peano_arithmetic_incompleteness}. Its proof is outside the scope of this monograph.

  \item Compare this to second-order Peano arithmetic, which is complete due to \cref{thm:def:hol_theory/categoricity_implies_completeness}.
\end{comments}

\begin{remark}\label{rem:thm:peano_arithmetic_incompleteness}
  \Fullref{thm:peano_arithmetic_incompleteness} is an important theorem with no known straightforward proof. Every sufficiently advanced \hyperref[con:logical_system]{logical system} subsumes Peano arithmetic, and hence suffers from the same incompleteness.

  Short after presenting his proof of \fullref{thm:fol_completeness}, G\"odel proved the incompleteness of \logic{PA} as \enquote{Satz VI} (\enquote{proposition VI}) in \cite{Gödel1931UnentscheidbareSätze}, in 1931. The gist of this proof was to encode first-order formulas and proof as natural numbers in order to craft a \hyperref[ex:self_reference]{self-referential} formula whose provability leads to a contradiction. Performing these constructions and showing that they are well-defined is a tedious task.

  Modern expositions of what is now known as \enquote{G\"odel's first incompleteness theorem} can be found, among other places, in \cite[thm. V]{Kleene2002Logic}, \cite[thm. 4.6.16]{Hinman2005Logic}, \cite[thm. 7.7.3]{VanDalen2004LogicAndStructure} and \cite[thm. 5.5.37]{Герасимов2014Вычислимость}. These authors state the theorem conditionally on the \hyperref[def:fol_theory/consistent]{consistency} of \logic{PA}, however, as noted in \cite[701]{GowersEtAl2008PrincetonCompanion}, this condition is redundant because any reasonable metatheory that allows proving incompleteness should also allow constructing a model of \logic{PA}.

  Later in the same paper, G\"odel also proved \enquote{Satz XI}, now known as \enquote{G\"odel's second incompleteness theorem} --- that a proof of consistency of \logic{PA}, encoded inside \logic{PA}, is not valid in the proof system encoded in \logic{PA}.

  Half a decade later, \tcite{in #2, #1}[thm. 1.3]{ParisHarrington1977Incompleteness} construct another formula unprovable in \logic{PA} --- a first-order formalization of a variant of \fullref{thm:ramseys_theorem}. Their theorem is seen as a more concrete alternative to G\"odel's first inconsistency theorem, however its proof relies on G\"odel's theorems.
\end{remark}

\paragraph{Order relations}

\begin{definition}\label{def:peano_arithmetic_ordering}
  We can introduce a new binary predicate \( \synleq \) to \hyperref[def:peano_arithmetic]{Peano arithmetic} using the \hyperref[def:fol_definition/predicate]{first-order definition}\fnote{Technically, we must take the \hyperref[def:fol_quantifier_closure]{universal closure} of \eqref{eq:def:peano_arithmetic_ordering/nonstrict}.}
  \begin{equation}\label{eq:def:peano_arithmetic_ordering/nonstrict}
    \synn \synleq \synm \syniff \qexists \syna (\synn \synplus \syna \syneq \synm).
  \end{equation}

  As we shall show in \cref{thm:natural_numbers_are_well_ordered}, this predicate acts as a \hyperref[def:well_ordered_set]{well-order}. We will show in \cref{thm:peano_arithmetic_order_compatibility} that the corresponding \hyperref[def:preordered_set/strict]{strict relation} can be defined via
  \begin{equation}\label{eq:def:peano_arithmetic_ordering/strict}
    \synn \synless \synm \syniff \qexists \syna (\synn \synplus \synS \syna \syneq \synm).
  \end{equation}
\end{definition}
\begin{comments}
  \item For the specific model of \logic{PA} based on the smallest inductive set \( \omega \), the latter relation \( \synless \) is interpreted as \( \in \); we discuss this in more detail in \cref{rem:ordinal_definition}.
\end{comments}

\begin{proposition}\label{thm:peano_arithmetic_order_rules}
  In \hyperref[def:peano_arithmetic]{Peano arithmetic} extended with \hyperref[def:peano_arithmetic_ordering]{order relations}, the following \hyperref[def:natural_deduction_rule]{natural deduction rules} are \hyperref[con:inference_rule_admissibility]{admissible}:
  \begin{thmenum}
    \thmitem{thm:peano_arithmetic_order_rules/refl} The nonstrict inequality is reflexive:
    \begin{equation*}\taglabel[\ensuremath{ \logic{PA}_{\logic{R}} }]{inf:thm:peano_arithmetic_order_rules/refl}
      \begin{prooftree}
        \infer0[\ref{inf:thm:peano_arithmetic_order_rules/refl}]{ \synn \synleq \synn }
      \end{prooftree}
    \end{equation*}

    \thmitem{thm:peano_arithmetic_order_rules/less_than_zero} No number is less than zero:
    \begin{equation*}\taglabel[\ensuremath{ \logic{PA}_{<0} }]{inf:thm:peano_arithmetic_order_rules/less_than_zero}
      \begin{prooftree}
        \hypo{ \synn \synless \syn0 }
        \infer1[\ref{inf:thm:peano_arithmetic_order_rules/less_than_zero}]{ \synbot }
      \end{prooftree}
    \end{equation*}

    \thmitem{thm:peano_arithmetic_order_rules/less_than_successor} Every number is smaller than its successor:
    \begin{equation*}\taglabel[\ensuremath{ \logic{PA}_{n<Sn} }]{inf:thm:peano_arithmetic_order_rules/less_than_successor}
      \begin{prooftree}
        \infer0[\ref{inf:thm:peano_arithmetic_order_rules/less_than_successor}]{ \synn \synless \synS \synn }
      \end{prooftree}
    \end{equation*}

    \thmitem{thm:peano_arithmetic_order_rules/nonstrict_successor} Successors of the upper bound are compatible with nonstrict inequalities:
    \begin{equation*}\taglabel[\ensuremath{ \logic{PA}_{m<Sn} }]{inf:thm:peano_arithmetic_order_rules/nonstrict_successor}
      \begin{prooftree}
        \hypo{ \synn \synleq \synm }
        \infer1[\ref{inf:thm:peano_arithmetic_order_rules/nonstrict_successor}]{ \synn \synleq \synS \synm }
      \end{prooftree}
    \end{equation*}

    \thmitem{thm:peano_arithmetic_order_rules/less_successor_to_leq} Strict inequality with a successor can be expressed via strict inequality:
    \begin{equation*}\taglabel[\ensuremath{ \logic{PA}_{< \to \leq} }]{inf:thm:peano_arithmetic_order_rules/less_successor_to_leq}
      \begin{prooftree}
        \hypo{ \synn \synless \synS \synm }
        \infer1[\ref{inf:thm:peano_arithmetic_order_rules/less_successor_to_leq}]{ \synn \synleq \synm }
      \end{prooftree}
    \end{equation*}

    \thmitem{thm:peano_arithmetic_order_rules/leq_to_less} Nonstrict inequality can be expressed via strict inequality:
    \begin{equation*}\taglabel[\ensuremath{ \logic{PA}_{\leq \to <} }]{inf:thm:peano_arithmetic_order_rules/leq_to_less}
      \begin{prooftree}
        \hypo{ \synn \synleq \synm }
        \infer1[\ref{inf:thm:peano_arithmetic_order_rules/leq_to_less}]{ \synn \synless \synm \synvee \synn \syneq \synm }
      \end{prooftree}
    \end{equation*}
  \end{thmenum}
\end{proposition}
\begin{proof}
  \SubProofOf{thm:peano_arithmetic_order_rules/refl}
  \begin{equation*}
    \begin{prooftree}
      \hypo{ [\eqref{eq:def:peano_arithmetic_ordering/nonstrict}]^{a_2} }
      \infer1[\ref{inf:def:fol_natural_deduction/forall/elim}*]{ n \synleq n \syniff \qexists a (n \synplus a \syneq n) }

      \hypo{ [\eqref{eq:def:peano_arithmetic/theory/PA4}]^{a_1} }
      \infer1[\ref{inf:def:fol_natural_deduction/forall/elim}]{ n \synplus \syn0 \syneq n }
      \infer1[\ref{inf:def:fol_natural_deduction/exists/intro}]{ \qexists a (n \synplus a \syneq n) }

      \infer2[\ref{inf:def:propositional_natural_deduction/iff/elim_right}]{ n \synleq n }
    \end{prooftree}
  \end{equation*}

  \SubProofOf{thm:peano_arithmetic_order_rules/less_than_zero}
  \footnotesize
  \begin{equation*}
    \begin{prooftree}[separation=1.2em]
      \hypo{ [n \synless \syn0]^u }
      \hypo{ [\eqref{eq:def:peano_arithmetic_ordering/strict}]^{a_1} }
      \infer1[\ref{inf:def:fol_natural_deduction/forall/elim}*]{ n \synless \syn0 \syniff \qexists a (n \synplus \synS a \syneq \syn0) }

      \infer2[\ref{inf:def:propositional_natural_deduction/iff/elim_right}]{ \qexists a (n \synplus \synS a \syneq \syn0) }
      \hypo{ [n \synplus \synS a \syneq \syn0]^v }

      \hypo{ [\eqref{eq:def:peano_arithmetic/theory/PA5}]^{a_2} }
      \infer1[\ref{inf:def:fol_natural_deduction/forall/elim}*]{ n \synplus \synS a \syneq \synS(n \synplus a) }
      \infer2[\ref{inf:def:fol_natural_deduction/equality/elim}]{ \synS(n \synplus a) \syneq \syn0 }
      \infer1[\ref{inf:def:fol_natural_deduction/exists/intro}]{ \qexists b (\synS b \syneq \syn0) }

      \hypo{ [\eqref{eq:def:peano_arithmetic/theory/PA1}]^{a_3} }
      \infer2[\ref{inf:def:propositional_natural_deduction/neg/elim}]{ \synbot }

      \infer[left label=\( v \)]2[\ref{inf:def:fol_natural_deduction/exists/elim}]{ \synbot }
    \end{prooftree}
  \end{equation*}
  \normalsize

  \SubProofOf{thm:peano_arithmetic_order_rules/less_than_successor}
  \begin{equation*}
    \begin{prooftree}
      \hypo{ [\eqref{eq:def:peano_arithmetic_ordering/strict}]^{a_2} }
      \infer1[\ref{inf:def:fol_natural_deduction/forall/elim}*]{ n \synless \synS n \syniff \qexists a (n \synplus \synS a \syneq \synS n) }

      \hypo{ [\eqref{eq:thm:peano_arithmetic_successor_and_addition}]^{a_1} }
      \infer1[\ref{inf:def:fol_natural_deduction/forall/elim}]{ \synS n \syneq n \synplus \synS \syn0 }
      \infer1[\ref{inf:thm:fol_natural_deduction_equality/symm}]{ n \synplus \synS \syn0 \syneq \synS n }
      \infer1[\ref{inf:def:fol_natural_deduction/exists/intro}]{ \qexists a (n \synplus \synS a \syneq \synS n) }

      \infer2[\ref{inf:def:propositional_natural_deduction/iff/elim_left}]{ n \synless \synS n }
    \end{prooftree}
  \end{equation*}

  \SubProofOf{thm:peano_arithmetic_order_rules/nonstrict_successor}
  \footnotesize
  \begin{equation*}
    \begin{prooftree}[separation=1.2em]
      \hypo{ [\eqref{eq:def:peano_arithmetic_ordering/strict}]^{a_1} }
      \infer1[\ref{inf:def:fol_natural_deduction/forall/elim}*]{ \synanon }

      \hypo{ [\eqref{eq:def:peano_arithmetic_ordering/strict}]^{a_1} }
      \infer1[\ref{inf:def:fol_natural_deduction/forall/elim}*]{ \synanon }

      \hypo{ [n \synleq m]^u }
      \infer2[\ref{inf:def:propositional_natural_deduction/iff/elim_right}]{ \qexists a (n \synplus \synS a \syneq m) }

      \hypo{ [n \synplus \synS a \syneq m]^v }
      \infer0[\ref{inf:def:fol_natural_deduction/equality/intro}]{ \synS m \syneq \synS m }
      \infer2[\ref{inf:def:fol_natural_deduction/equality/elim}]{ \synS(n \synplus \synS a) \syneq \synS m }

      \hypo{ [\eqref{eq:def:peano_arithmetic/theory/PA5}]^{a_2} }
      \infer1[\ref{inf:def:fol_natural_deduction/forall/elim}*]{ \synanon }
      \infer2[\ref{inf:def:fol_natural_deduction/equality/elim}]{ n \synplus \synS \synS a \syneq \synS m }
      \infer1[\ref{inf:def:fol_natural_deduction/exists/intro}]{ \qexists b (n \synplus \synS b \syneq \synS m) }

      \infer[left label=\( v \)]2[\ref{inf:def:fol_natural_deduction/exists/elim}]{ \qexists b (n \synplus \synS b \syneq \synS m) }

      \infer2[\ref{inf:def:propositional_natural_deduction/iff/elim_left}]{ n \synless \synS m }
    \end{prooftree}
  \end{equation*}
  \normalsize

  \SubProofOf{thm:peano_arithmetic_order_rules/less_successor_to_leq}
  \footnotesize
  \begin{equation*}
    \begin{prooftree}[separation=1.2em]
      \hypo{ [\eqref{eq:def:peano_arithmetic_ordering/nonstrict}]^{a_4} }
      \infer1[\ref{inf:def:fol_natural_deduction/forall/elim}*]{ \synanon }

      \hypo{ [m \synless \synS n]^u }

      \hypo{ [\eqref{eq:def:peano_arithmetic_ordering/strict}]^{a_1} }
      \infer1[\ref{inf:def:fol_natural_deduction/forall/elim}*]{ \synanon }
      \infer2[\ref{inf:def:propositional_natural_deduction/iff/elim_right}]{ \qexists a (m \synplus \synS a \syneq \synS n) }

      \hypo{ [m \synplus \synS a \syneq \synS n]^v }

      \hypo{ [\eqref{eq:def:peano_arithmetic/theory/PA5}]^{a_2} }
      \infer1[\ref{inf:def:fol_natural_deduction/forall/elim}]{ \synanon }
      \infer2[\ref{inf:def:fol_natural_deduction/equality/elim}]{ \synS(m \synplus a) \syneq \synS n }

      \hypo{ [\eqref{eq:def:peano_arithmetic/theory/PA2}]^{a_3} }
      \infer1[\ref{inf:def:fol_natural_deduction/forall/elim}*]{ \synanon }
      \infer2[\ref{inf:def:propositional_natural_deduction/imp/elim}]{ m \synplus a \syneq n }
      \infer1[\ref{inf:def:fol_natural_deduction/exists/intro}]{ \qexists a (m \synplus a \syneq n) }
      \infer[left label=\( v \)]2[\ref{inf:def:fol_natural_deduction/exists/elim}]{ \qexists a (m \synplus a \syneq n) }

      \infer2[\ref{inf:def:propositional_natural_deduction/iff/elim_right}]{ m \synleq n }
    \end{prooftree}
  \end{equation*}
  \normalsize

  \SubProofOf{thm:peano_arithmetic_order_rules/leq_to_less}
  \begin{equation*}
    \begin{prooftree}
      \hypo{ [m \synleq n]^u }

      \hypo{ [\eqref{eq:def:peano_arithmetic_ordering/nonstrict}]^{a_1} }
      \infer1[\ref{inf:def:fol_natural_deduction/forall/elim}*]{ \synanon }
      \infer2[\ref{inf:def:propositional_natural_deduction/iff/elim_right}]{ \qexists a (m \synplus a \syneq n) }

      \hypo{ [\eqref{eq:thm:peano_arithmetic_predecessor_existence}]^{a_2} }
      \infer1[\ref{inf:def:fol_natural_deduction/forall/elim}]{ \synanon }

      \hypo{ \substack{ [m \synplus a \syneq n]^v \\ [a \syneq \syn0]^a } }
      \ellipsis {\( L \)} { m \syneq n }
      \infer1[\ref{inf:def:propositional_natural_deduction/or/intro_right}]{ \synanon }

      \hypo{ \substack{ [m \synplus a \syneq n]^v \\ [\qexists b (a \syneq \synS b)]^b } }
      \ellipsis {\( R \)} { m \synless n }
      \infer1[\ref{inf:def:propositional_natural_deduction/or/intro_left}]{ \synanon }

      \infer[left label={\( a, b \)}]3[\ref{inf:def:propositional_natural_deduction/or/elim}]{ (m \synless n) \synvee (m \syneq n) }

      \infer[left label=\( v \)]2[\ref{inf:def:fol_natural_deduction/exists/elim}]{ (m \synless n) \synvee (m \syneq n) }
    \end{prooftree}
  \end{equation*}
  where the subtree \( L \) is
  \begin{equation*}
    \begin{prooftree}
      \hypo{ [m \synplus a \syneq n]^v }
      \hypo{ [a \syneq \syn0]^a }
      \infer2[\ref{inf:def:fol_natural_deduction/equality/elim}]{ m \synplus \syn0 \syneq n }

      \hypo{ [\eqref{eq:def:peano_arithmetic/theory/PA4}]^{a_3} }
      \infer1[\ref{inf:def:fol_natural_deduction/forall/elim}*]{ \synanon }
      \infer2[\ref{inf:def:fol_natural_deduction/equality/elim}]{ m \syneq n }
    \end{prooftree}
  \end{equation*}
  and \( R \) is
  \begin{equation*}
    \begin{prooftree}
      \hypo{ [\qexists b (a \syneq \synS b)]^b }

      \hypo{ [\eqref{eq:def:peano_arithmetic_ordering/strict}]^{a_4} }
      \infer1[\ref{inf:def:fol_natural_deduction/forall/elim}*]{ \synanon }

      \hypo{ [m \synplus a \syneq n]^v }
      \hypo{ [a \syneq \synS c]^c }
      \infer2[\ref{inf:def:fol_natural_deduction/equality/elim}]{ m \synplus \synS c \syneq n }
      \infer1[\ref{inf:def:fol_natural_deduction/exists/intro}]{ \qexists c (m \synplus \synS c \syneq n) }

      \infer2[\ref{inf:def:propositional_natural_deduction/iff/elim_left}]{ m \synless n }

      \infer[left label=\( c \)]2[\ref{inf:def:fol_natural_deduction/exists/elim}]{ m \synless n }
    \end{prooftree}
  \end{equation*}
\end{proof}

\begin{proposition}\label{thm:peano_arithmetic_order_compatibility}
  In \hyperref[def:peano_arithmetic]{Peano arithmetic}, the inequality predicates from \cref{def:peano_arithmetic_ordering} are indeed related via \eqref{eq:def:preordered_set/compatibility_nonstrict}.
\end{proposition}
\begin{proof}
  Our first step is to prove by induction the following (we keep \( \synn \) free):
  \begin{equation}\label{eq:thm:peano_arithmetic_order_compatibility/proof/aux}
    \qforall \synm \varphi[\synx \mapsto \synm],
  \end{equation}
  where
  \begin{equation*}
    \varphi \coloneqq (\synn \synless \synx \synvee \synn \syneq \synx) \synimplies \synn \synleq \synx
  \end{equation*}

  \SubProofOf{eq:thm:peano_arithmetic_order_compatibility/proof/aux}
  \begin{equation*}
    \begin{prooftree}
      \hypo{ }
      \ellipsis {\( L \)} { \varphi[\synx \mapsto \syn0] }

      \hypo{ }
      \ellipsis {\( R \)} { \varphi[\synx \mapsto \synS \synm] }
      \infer1[\ref{inf:def:propositional_natural_deduction/imp/intro}]{ \varphi[\synx \mapsto \synm] \synimplies \varphi[\synx \mapsto \synS \synm] }
      \infer1[\ref{inf:def:fol_natural_deduction/forall/intro}]{ \qforall \synm (\varphi[\synx \mapsto \synm] \synimplies \varphi[\synx \mapsto \synS \synm]) }

      \infer2[\ref{inf:thm:peano_arithmetic_induction_rule}]{ \qforall \synm \syn\varphi[\synx \mapsto \synm] }
    \end{prooftree}
  \end{equation*}
  where the subtree \( L \) is
  \begin{equation*}
    \begin{prooftree}
      \hypo{ [\synn \synless \syn0 \synvee \synn \syneq \syn0]^b }

      \hypo{ [\synn \synless \syn0]^c }
      \infer1[\ref{inf:thm:peano_arithmetic_order_rules/less_than_zero}]{ \synbot }
      \infer1[\ref{inf:def:propositional_natural_deduction/bot/efq}]{ \synn \synleq \syn0 }

      \hypo{ [\synn \syneq \syn0]^d }
      \infer0[\ref{inf:thm:peano_arithmetic_order_rules/refl}]{ \synn \synleq \synn }
      \infer2[\ref{inf:def:fol_natural_deduction/equality/elim}]{ \synn \synleq \syn0 }

      \infer[left label={\( c, d \)}]3[\ref{inf:def:propositional_natural_deduction/or/elim}]{ \synn \synleq \syn0 }
      \infer[left label={\( b \)}]1[\ref{inf:def:propositional_natural_deduction/imp/intro}]{ \varphi[\synx \mapsto \syn0] }
    \end{prooftree}
  \end{equation*}
  and \( R \) is
  \begin{equation*}
    \begin{prooftree}
      \hypo{ [\synn \synless \synS \synm \synvee \synn \syneq \synS \synm]^e }

      \hypo{ [\synn \synless \synS \synm]^f }
      \infer1[\ref{inf:thm:peano_arithmetic_order_rules/less_successor_to_leq}]{ \synn \synleq \synm }
      \infer1[\ref{thm:peano_arithmetic_order_rules/nonstrict_successor}]{ \synn \synleq \synS \synm }

      \hypo{ [\synn \syneq \synS \synm]^g }
      \infer0[\ref{inf:thm:peano_arithmetic_order_rules/refl}]{ \synn \synleq \synn }
      \infer2[\ref{inf:def:fol_natural_deduction/equality/elim}]{ \synn \synleq \synS \synm }

      \infer[left label={\( f, g \)}]3[\ref{inf:def:propositional_natural_deduction/or/elim}]{ \synn \synleq \synS \synm }
      \infer[left label={\( e \)}]1[\ref{inf:def:propositional_natural_deduction/imp/intro}]{ \varphi[\synx \mapsto \synS \synm] }
    \end{prooftree}
  \end{equation*}

  \SubProofOf{eq:def:preordered_set/compatibility_nonstrict}
  \begin{equation*}
    \begin{prooftree}
      \hypo{ [\synn \synleq \synm]^u }
      \infer1[\ref{inf:thm:peano_arithmetic_order_rules/leq_to_less}]{ \synn \synless \synm \synvee \synn \syneq \synm }

      \hypo{ [\eqref{eq:thm:peano_arithmetic_order_compatibility/proof/aux}]^{a_1} }
      \infer1[\ref{inf:def:fol_natural_deduction/forall/elim}*]{ \synanon }

      \hypo{ [\synn \synless \synm \synvee \synn \syneq \synm]^v }
      \infer2[\ref{inf:def:propositional_natural_deduction/imp/elim}]{ \synn \synleq \synm }

      \infer[left label={\( u, v \)}]2[\ref{inf:def:propositional_natural_deduction/iff/intro}]{ (\synn \synleq \synm) \syniff (\synn \synless \synm \synvee \synn \syneq \synm) },
      \infer1[\ref{inf:def:fol_natural_deduction/forall/intro}*]{ \eqref{eq:def:preordered_set/compatibility_nonstrict} }
    \end{prooftree}
  \end{equation*}
\end{proof}

\paragraph{Strong induction}

\begin{theorem}[Strong induction on natural numbers]\label{thm:strong_induction_on_natural_numbers}
  In \hyperref[def:peano_arithmetic]{Peano arithmetic}, we can \hyperref[def:fol_natural_deduction_consequence]{formally derive} the following induction principle
  \begin{equation}\label{eq:thm:strong_induction_on_natural_numbers}\tag{\ensuremath{ \logic{PA3}^{\logic{S}} }}
    \qforall \synn \parens[\big]
      {
        \overbrace
          {
            \underbrace{ \qforall \synm \parens[\big]{ \synm \synless \synn \synimplies \syn\varphi[\synx \synsubst \synm] } }_{\mathclap{\substack{\T{inductive} \\ \T{hypothesis}}}}
            \synimplies
            \underbrace{ \syn\varphi[\synx \synsubst \synn] }_{\mathclap{\substack{\T{inductive step} \\ \T{conclusion}}}}
          }^{\T{inductive step}}
      }
    \synimplies
    \underbrace{ \qforall \synn \syn\varphi[\synx \synsubst \synn] }_{\T{conclusion}}.
  \end{equation}

  The corresponding \hyperref[def:second_order_logic]{second-order axiom} is
  \begin{equation}\label{eq:thm:strong_induction_on_natural_numbers/sol}\tag{\ensuremath{ \logic{PA3}_2^{\logic{S}} }}
    \qforall {\synp^{\syn\iota \synimplies \syn\omicron}}
    \parens[\big]
      {
        \qforall {\synn^{\syn\iota}} \parens[\big]
          {
            \qforall {\synm^{\syn\iota}} (\synm \synless \synn \synimplies \synp \synm)
            \synimplies
            \synp \synn
          }
        \synimplies
        \qforall {\synn^{\syn\iota}} \synp \synn
      }.
  \end{equation}

  \Fullref{thm:induction_on_natural_numbers} can be restated as follows:
  \begin{displayquote}
    In order for \( p(n) \) to hold \hi{for every} \( n \) in \( N \), it is sufficient for \( p(n) \) to hold \hi{whenever} \( p(m) \) holds for every \( m < n \).
  \end{displayquote}
\end{theorem}
\begin{comments}
  \item Strong induction can be regarded as special case of \fullref{thm:bounded_transfinite_induction}. See \cref{con:induction} for a general discussion of induction and recursion.

  \item \Cref{thm:omega_plus_one_model_of_robinson_arithmetic} shows that strong induction is by itself not sufficient to derive the usual induction principle, however \cref{thm:strong_induction_implies_induction} gives a sufficient condition.

  \item We need multiple instances of \eqref{eq:def:peano_arithmetic/theory/PA3} to derive one instance of \eqref{eq:thm:strong_induction_on_natural_numbers}.

  \item This theorem is stated for \( \omega \) in \cite[87]{Enderton1977SetTheory}.
\end{comments}
\begin{proof}
  Fix a formula \( \varphi \). Without loss of generality, suppose that \( x \) is the only free variable of \( \varphi \); otherwise we will have to apply \ref{inf:def:fol_natural_deduction/forall/intro} and \ref{inf:def:fol_natural_deduction/forall/elim} the appropriate number of times.

  For brevity, we will use the (metalinguistic) abbreviations
  \begin{align*}
    \varphi_n &\coloneqq \varphi[x \mapsto n], \\
    \psi_n    &\coloneqq \qforall m (m \synless n \synimplies \varphi_m).
  \end{align*}

  We will again need an intermediate step --- the following variation of \eqref{eq:thm:strong_induction_on_natural_numbers}:
  \begin{equation}\label{eq:thm:strong_induction_on_natural_numbers/proof/aux}
    \qforall n (\psi_n \synimplies \varphi_n) \synimplies \qforall n \psi_n
  \end{equation}

  \SubProofOf{eq:thm:strong_induction_on_natural_numbers/proof/aux}
  \begin{equation*}
    \begin{prooftree}
      \hypo{ [m \synless \syn0]^u }
      \infer1[\ref{inf:thm:peano_arithmetic_order_rules/less_than_zero}]{ \synbot }
      \infer1[\ref{inf:def:propositional_natural_deduction/bot/efq}]{ \varphi_m }
      \infer[left label={\( v \)}]1[\ref{inf:def:propositional_natural_deduction/imp/intro}]{ m \synless \syn0 \synimplies \varphi_m }
      \infer1[\ref{inf:def:fol_natural_deduction/forall/intro}]{ \psi_0 }

      \hypo{ \substack{ [\qforall n (\psi_n \synimplies \varphi_n)]^v \\ [\psi_n]^w } }
      \ellipsis {\( R \)} { \psi_{\synS n} }
      \infer[left label={\( w \)}]1[\ref{inf:def:propositional_natural_deduction/imp/intro}]{ \psi_n \synimplies \psi_{\synS n} }
      \infer1[\ref{inf:def:fol_natural_deduction/forall/intro}]{ \qforall n (\psi_n \synimplies \psi_{\synS n}) }

      \infer2[\ref{inf:thm:peano_arithmetic_induction_rule}]{ \qforall n \psi_n }
      \infer[left label={\( v \)}]1[\ref{inf:def:propositional_natural_deduction/imp/intro}]{ \qforall n (\psi_n \synimplies \varphi_n) \synimplies \qforall n \psi_n }
    \end{prooftree}
  \end{equation*}
  where the subtree \( R \) is
  \scriptsize
  \begin{equation*}
    \begin{prooftree}[separation=1.5em]
      \hypo{ [m \synless \synS n]^b }
      \infer1[\ref{inf:thm:peano_arithmetic_order_rules/less_successor_to_leq}]{ m \synleq n }
      \infer1[\ref{inf:thm:peano_arithmetic_order_rules/leq_to_less}]{ m \synless n \synvee m \syneq n }

      \hypo{ [\psi_n]^w }
      \infer1[\ref{inf:def:fol_natural_deduction/forall/elim}]{ \synanon }

      \hypo{ [m \synless n]^c }
      \infer2[\ref{inf:def:propositional_natural_deduction/imp/elim}]{ \varphi_m }

      \hypo{ [\synanon]^v }
      \infer1[\ref{inf:def:fol_natural_deduction/forall/elim}]{ \psi_n \synimplies \varphi_n }
      \hypo{ [\psi_n]^w }

      \infer2[\ref{inf:def:propositional_natural_deduction/imp/elim}]{ \varphi_n }

      \hypo{ [m \syneq n]^d }
      \infer2[\ref{inf:def:propositional_natural_deduction/imp/elim}]{ \varphi_m }

      \infer[left label={\( c, d \)}]3[\ref{inf:def:propositional_natural_deduction/or/elim}]{ \varphi_m }
      \infer[left label={\( b \)}]1[\ref{inf:def:propositional_natural_deduction/imp/intro}]{ \psi_{\synS n} }
    \end{prooftree}
  \end{equation*}
  \normalsize

  \SubProofOf{eq:thm:strong_induction_on_natural_numbers}
  \begin{equation*}
    \begin{prooftree}
      \hypo{ [\eqref{eq:thm:strong_induction_on_natural_numbers/proof/aux}]^{a_1} }
      \infer1[\ref{inf:def:fol_natural_deduction/forall/elim}]{ \synanon }

      \hypo{ [\qforall n (\psi_n \synimplies \varphi_n)]^u }
      \infer2[\ref{inf:def:propositional_natural_deduction/imp/elim}]{ \qforall n \psi_n }
      \infer1[\ref{inf:def:fol_natural_deduction/forall/elim}]{ \psi_{\synS n} }
      \infer1[\ref{inf:def:fol_natural_deduction/forall/elim}]{ n \synless \synS n \synimplies \varphi_n }

      \infer0[\ref{inf:thm:peano_arithmetic_order_rules/less_than_successor}]{ n \synless \synS n }
      \infer2[\ref{inf:def:propositional_natural_deduction/imp/elim}]{ \varphi_n }
      \infer1[\ref{inf:def:fol_natural_deduction/forall/intro}]{ \qforall \synn \varphi_n }
      \infer[left label={\( u \)}]1[\ref{inf:def:propositional_natural_deduction/imp/intro}]{ \eqref{eq:thm:strong_induction_on_natural_numbers} }
    \end{prooftree}
  \end{equation*}
\end{proof}

\begin{proposition}\label{thm:omega_plus_one_model_of_robinson_arithmetic}\mimprovised
  Consider the infinite \hyperref[def:ordinal]{ordinal} \( \omega + 1 \) and define
  \begin{align*}
    \alpha \oplus \beta &\coloneqq \begin{cases}
      \alpha + \beta, &\alpha < \omega \T{and} \beta < \omega, \\
      \omega,         &\T{otherwise.}
    \end{cases} \\
    \alpha \odot \beta &\coloneqq \begin{cases}
      \alpha \cdot \beta, &\alpha < \omega \T{and} \beta < \omega, \\
      \omega,         &\T{otherwise.}
    \end{cases} \\
    S(\alpha) &\coloneqq \alpha \oplus 1.
  \end{align*}

  With these operations, \( \omega + 1 \) is a \hyperref[def:fol_semantics/model]{model} of \hyperref[def:peano_arithmetic]{Robinson arithmetic} with the strong induction schema \eqref{eq:thm:strong_induction_on_natural_numbers}.
\end{proposition}
\begin{comments}
  \item This is not a model of Peano arithmetic because \( \omega \) is a limit ordinal and thus cannot be reached from \( 0 \) by applying \( S \) iteratively.
\end{comments}
\begin{proof}
  The Robinson arithmetic axioms are validated by \fullref{thm:omega_is_model_of_pa} for finite ordinals and are trivially satisfied for \( \omega \).

  \Cref{thm:bounded_transfinite_induction} shows that strong induction holds in \( \omega + 1 \).
\end{proof}

\begin{proposition}\label{thm:strong_induction_implies_induction}
  In \hyperref[def:peano_arithmetic]{Robinson arithmetic}, the strong induction schema \eqref{eq:thm:strong_induction_on_natural_numbers} together with the predecessor existence formula \eqref{eq:thm:peano_arithmetic_predecessor_existence} together allow deriving any instance of the usual induction principle \eqref{eq:def:peano_arithmetic/theory/PA3}.
\end{proposition}
\begin{proof}
  As in \fullref{thm:strong_induction_on_natural_numbers}, fix a formula \( \varphi \), assumed to have only one free variable, \( x \). We will reuse the notational abbreviation \( \varphi_n \coloneqq \varphi[x \mapsto n] \).

  Then our desired derivation is
  \begin{equation*}
    \begin{prooftree}
      \hypo{ [\eqref{eq:thm:strong_induction_on_natural_numbers}]^{a_1} }

      \hypo{ [\eqref{eq:thm:peano_arithmetic_predecessor_existence}]^{a_2} }
      \infer1[\ref{inf:def:fol_natural_deduction/forall/elim}]{ n \syneq \syn0 \synvee \qexists a (n \syneq \synS a) }

      \hypo{ \substack{ [\varphi_0 \synwedge \qforall n (\varphi_n \synimplies \varphi_{\synS n})]^u \\ [\qforall m (m \synless n \synimplies \varphi_m)]^v \\ [n \syneq \syn0]^b } }
      \ellipsis {\( L \)} { \varphi_n }

      \hypo{ \substack{ [\varphi_0 \synwedge \qforall n (\varphi_n \synimplies \varphi_{\synS n})]^u \\ [\qforall m (m \synless n \synimplies \varphi_m)]^v \\ [\qexists a (n \syneq \synS a)]^c } }
      \ellipsis {\( R \)} { \varphi_n }

      \infer[left label={\( b, c \)}]3[\ref{inf:def:propositional_natural_deduction/or/elim}]{ \varphi_n }
      \infer[left label={\( v \)}]1[\ref{inf:def:propositional_natural_deduction/imp/intro}]{ \qforall m (m \synless n \synimplies \varphi_m) \synimplies \varphi_n }
      \infer1[\ref{inf:def:fol_natural_deduction/forall/intro}]{ \qforall n (\qforall m (m \synless n \synimplies \varphi_m) \synimplies \varphi_n) }

      \infer2[\ref{inf:def:propositional_natural_deduction/imp/elim}]{ \qforall n \varphi_n }
      \infer[left label={\( u \)}]1[\ref{inf:def:propositional_natural_deduction/imp/intro}]{ \eqref{eq:def:peano_arithmetic/theory/PA3} }
    \end{prooftree}
  \end{equation*}
  where the subtree \( L \) is
  \begin{equation*}
    \begin{prooftree}
      \hypo{ [n \syneq \syn0]^b }
      \hypo{ [\varphi_0 \synwedge \qforall n (\varphi_n \synimplies \varphi_{\synS n})]^u }
      \infer1[\ref{inf:def:propositional_natural_deduction/and/elim_left}]{ \varphi_{\syn0} }
      \infer2[\ref{inf:def:fol_natural_deduction/equality/elim}]{ \varphi_n }
    \end{prooftree}
  \end{equation*}
  and \( R \) is
  \footnotesize
  \begin{equation*}
    \begin{prooftree}[separation=1em]
      \hypo{ [\qexists a (n \syneq \synS a)]^c }

      \hypo{ [\synanon]^v }
      \infer1[\ref{inf:def:fol_natural_deduction/forall/elim}]{ b \synless n \synimplies \varphi_b }

      \hypo{ [n \syneq \synS b]^d }
      \infer0[\ref{inf:thm:peano_arithmetic_order_rules/less_than_successor}]{ b \synless \synS b }
      \infer2[\ref{inf:def:fol_natural_deduction/equality/elim}]{ b \synless n }

      \infer2[\ref{inf:def:propositional_natural_deduction/imp/elim}]{ \varphi_b }

      \hypo{ [\synanon]^u }
      \infer1[\ref{inf:def:propositional_natural_deduction/and/elim_right}]{ \qforall n (\varphi_n \synimplies \varphi_{\synS n}) }
      \infer1[\ref{inf:def:fol_natural_deduction/forall/elim}]{ \varphi_b \synimplies \varphi_{\synS b} }

      \hypo{ [n \syneq \synS b]^d }
      \infer2[\ref{inf:def:fol_natural_deduction/equality/elim}]{ \varphi_b \synimplies \varphi_n }

      \infer2[\ref{inf:def:propositional_natural_deduction/imp/elim}]{ \varphi_n }
      \infer[left label=\( d \)]2[\ref{inf:def:fol_natural_deduction/exists/elim}]{ \varphi_n }
    \end{prooftree}
  \end{equation*}
  \normalsize
\end{proof}
